% =====================================================================
%  Avance 1 — Control de acceso estudiantil por código QR (TESCHA)
%  Plantilla APA 7.ª edición (clase apa7)
%  Compilar con:  latexmk -pdf trabajo.tex   (usa biber automáticamente)
% =====================================================================
\documentclass[stu,11pt,floatsintext,biblatex]{apa7}

% Fuente sans serif a 11 pt (APA permite sans serif de 11 pt, como Arial)
\usepackage{tgheros}
\renewcommand{\familydefault}{\sfdefault}

\usepackage[spanish,es-tabla]{babel} % es-tabla: "Tabla" en vez de "Cuadro"
\usepackage{csquotes}                % requerido por biblatex-apa
\addbibresource{citas.bib}

\usepackage{tikz}
\usetikzlibrary{arrows.meta,positioning,shapes.geometric}

% La clase apa7 deja el texto en bandera (raggedright) tras \maketitle;
% aquí se justifica, como pide el profesor.
\usepackage{ragged2e}

% Logos arriba de la portada (los pide la escuela, no APA).
\newcommand{\hrm}{\noindent
  \includegraphics[height=1.6cm]{img/tescha_logo.png}\hfill
  \includegraphics[height=1.6cm]{img/logo_sitemas.png}\par}

% Índice después de la portada (lo pide la escuela; APA puro no lo usa).
\setcounter{tocdepth}{2}
\makeatletter
\patchcmd{\maketitle}{\@ifundefined{def@donotrepeattitle}}
  {\tableofcontents\newpage\@ifundefined{def@donotrepeattitle}}{}{}
\patchcmd{\maketitle}{\section*{\protect\normalfont\textbf{\@title}}}
  {\section*{\protect\normalfont\textbf{\@title}}\addcontentsline{toc}{section}{\@title}}{}{}
\makeatother
\AtBeginDocument{\defbibheading{bibliography}[\refname]{%
  \clearpage\section*{\normalfont\textbf{#1}}\addcontentsline{toc}{section}{#1}}}
\newcommand{\apendiceindice}[2]{\addcontentsline{toc}{section}{\appendixname~#1: #2}}

\fancypagestyle{titlepage}{\fancyhf{}\renewcommand{\headrulewidth}{0pt}}

\renewcommand{\figurenote}[1]{\par\noindent\textit{\notelabel.} #1\par}
\renewcommand{\tablenote}[1]{\par\smallskip\figurenote{#1}}

\hypersetup{hidelinks}
\setlength{\headheight}{15.4pt}

\DeclareDelimFormat[parencite]{finalnamedelim}{\addspace\bibstring{and}\space}
\DeclareDelimFormat[bib,biblist]{finalnamedelim}{%
  \ifthenelse{\value{listcount}>\maxprtauth}{}{\addspace\bibstring{and}\space}}
\DefineBibliographyExtras{spanish}{\protected\def\bibrangedash{\textendash}}

\AtBeginDocument{%
  \renewcommand{\keywordname}{Palabras clave}%
  \renewcommand{\acksname}{Nota del autor}%
  \renewcommand{\notelabel}{Nota}%
}

% ----------------------------- Datos ---------------------------------
\title{Escanea y Entra: Prototipo de Acceso QR para el TESCHA}
\shorttitle{ACCESO QR TESCHA}

\authorsnames{Juan Esteban Anzures Campos, Daniel Arellano Palacios, Ian Alexis Flores Martínez, Miguel Ángel Jurado Delgadillo, Juan Fernando Reyes Arguello, Bryan Uriel Villegas Raymundo}
\authorsaffiliations{{Ingeniería en Sistemas Computacionales, Tecnológico de Estudios Superiores de Chalco}}

% Modo stu
\course{Investigación (grupo 4751, turno vespertino)}
\professor{Dra. en C. C. Giselle Nayeli Morales Rosales}
\duedate{30 de septiembre de 2026}

\makeatletter
\@ifundefined{def@stu}{}{\let\@shorttitle\@undefined \let\@acks\@undefined}
\makeatother

% =====================================================================
% Casilla para opciones de la encuesta
\newcommand{\op}{\framebox[0.8em]{\rule{0pt}{0.55em}}\,}
\newcommand{\linea}[1][4cm]{\rule{#1}{0.4pt}}

\begin{document}
\maketitle
\justifying

Este avance corresponde al proyecto integrador cuyo título provisional es
\textit{Escanea y Entra: Prototipo de Acceso QR para el TESCHA}. Siguiendo los
pasos para construir el inicio de un proyecto de investigación, en este
documento se define el tema de interés, se delimita el proyecto, se plantea
el problema de investigación —su situación, su importancia, lo que
desconocemos de ella y su viabilidad—, se formula la pregunta de
investigación y se explica cómo se investigará. Los instrumentos con los que
se levantarán los datos del plantel se incluyen en los apéndices.


\clearpage
\section{Tema}
El tema de interés general de este proyecto es el uso de tecnologías de
identificación automática —como el código QR, las tarjetas de proximidad
(NFC) o el reconocimiento facial— para el control de acceso de personas en
instituciones educativas. Estas tecnologías buscan sustituir la revisión
visual de credenciales físicas, un método que depende por completo del
criterio y la atención del personal de vigilancia y cuya insatisfacción ya se
ha documentado en otros planteles \parencite{torres2019qr}. En educación
superior se han propuesto sistemas basados en código QR, combinado con otras
tecnologías, para agilizar el registro de los estudiantes
\parencite{cano2026sistema}, y el uso de códigos QR para el control de acceso
físico en general ya se ha planteado en la literatura
\parencite{kao2011physical}.


\section{Delimitación}
Dentro de ese tema, el proyecto se delimita a la validación de identidad en
un acceso peatonal del Tecnológico de Estudios Superiores de Chalco (TESCHA),
en Chalco, Estado de México. Se elige un solo acceso, en el turno vespertino,
como punto piloto tanto para la observación inicial como para la prueba del
prototipo (Apéndices~A, B y~C), y se opta por el código QR dinámico como
tecnología de identificación, según se justifica en ``La Viabilidad'', más
adelante.

Concretamente, se diseñará y construirá un prototipo de torniquete
controlado por un microcontrolador ESP32 con cámara (ESP32-CAM), capaz de
leer un código QR dinámico asignado a cada estudiante, validarlo contra un
servidor y registrar cada ingreso en una base de datos consultable desde un
panel web (Figura~\ref{fig:arquitectura}). El código se regenera cada pocos
segundos y se firma con una contraseña de un solo uso basada en HMAC, un
esquema que resiste la interceptación, el reenvío y la modificación del
mensaje \parencite{subpratatsavee2015hmac}; así, una captura de pantalla
compartida deja de ser válida al poco tiempo. Esta estrategia de regeneración
periódica ya se ha probado en el control de asistencia de una universidad del
Reino Unido, donde un código QR renovado cada 20 segundos impidió que los
estudiantes reenviaran capturas de pantalla entre sí
\parencite{nwabuwe2023dynamicqr}; otros trabajos han combinado de forma
similar el QR con una contraseña de un solo uso para el acceso a laboratorios
universitarios \parencite{chaabani2023labqr}.

\begin{figure}[tbp]
	\caption{Arquitectura Propuesta del Sistema de Acceso por Código QR}
	\label{fig:arquitectura}
	\centering
	\begin{tikzpicture}[
		node distance=1.1cm and 1.6cm,
		box/.style={draw, rounded corners, align=center, minimum height=1cm,
				minimum width=2.6cm, font=\small},
		arr/.style={-{Latex[length=2mm]}, thick}
		]
		\node[box] (alumno) {QR dinámico\\(celular del alumno)};
		\node[box, right=of alumno] (cam) {ESP32-CAM\\(captura y decodifica)};
		\node[box, right=of cam] (esp) {ESP32\\(controla torniquete)};
		\node[box, below=of esp] (servidor) {Servidor y BD\\(valida el token)};
		\node[box, left=of servidor] (panel) {Panel web\\administrativo};
		\draw[arr] (alumno) -- (cam) node[midway, above, font=\scriptsize]{escaneo};
		\draw[arr] (cam) -- (esp) node[midway, above, font=\scriptsize]{token};
		\draw[arr] (esp) -- (servidor) node[midway, right, font=\scriptsize]{Wi-Fi/HTTPS};
		\draw[arr] (servidor) -- (panel) node[midway, below, font=\scriptsize]{consulta};
	\end{tikzpicture}
	\figurenote{Elaboración propia. El servidor responde si el token es válido
		y el ESP32 libera el torniquete; cada intento, aceptado o rechazado, queda
		registrado en la base de datos.}
\end{figure}

La delimitación del proyecto en esta etapa es la siguiente:
\begin{itemize}
	\item \textbf{Incluye:} un acceso piloto, el registro de entradas de
	      estudiantes, la generación y validación de QR dinámicos, la base de datos de
	      ingresos y un panel web de consulta.
	\item \textbf{No incluye:} reconocimiento facial, tarjetas NFC, pase de lista
	      por materia, registro de salidas ni control de visitantes. Los estudiantes
	      que no puedan mostrar su código seguirán el procedimiento manual actual.
\end{itemize}

Cabe reconocer una limitación: un QR dinámico impide reutilizar capturas de
pantalla, pero no que un estudiante preste su celular a otra persona. Por esa
razón, el proyecto no pretende eliminar la suplantación, sino ayudar a
aliviar los estragos que genera el sistema actual.


\section{Problema de Investigación}
Para plantear el problema que motiva este proyecto se describe primero la
situación actual, después su importancia, en seguida lo que aún desconocemos
de ella y, por último, la viabilidad de estudiarla e intervenirla.

\subsection{La Situación}
El ingreso de los estudiantes al Tecnológico de Estudios Superiores de Chalco
(TESCHA) se controla mediante la revisión visual de credenciales físicas por
parte del personal de vigilancia. Este método es lento en las horas de mayor
afluencia, no deja un registro digital de quién entra ni a qué hora y no
permite comprobar que quien presenta la credencial sea su titular. La
insatisfacción con los registros manuales se ha documentado en otros
contextos: en una escuela de educación básica de Perú, el 87.5\,\% de los 40
docentes y administrativos encuestados no se sentía cómodo con el registro
manual de asistencia \parencite{torres2019qr}. En el TESCHA, la magnitud de
estas limitaciones se medirá con los instrumentos de los Apéndices~A, B y~C.

El problema identificado es que el acceso de los estudiantes al TESCHA carece
de un mecanismo de validación automatizado. Como consecuencia, se forman
filas en los minutos previos al inicio de clases, una credencial prestada o
ajena puede superar la revisión visual y la institución no cuenta con un
registro confiable de ingresos que sirva para tomar decisiones.

Conviene distinguir este problema del pase de lista dentro del salón (la
práctica conocida como ``firmar por un compañero''). Ambos comparten la misma
raíz, una identificación deficiente, pero un registro en la entrada solo
demuestra que el estudiante ingresó al plantel, no que asistió a una clase
específica. Por ello, el control de asistencia por materia queda fuera del
alcance de este proyecto.

Esta situación se presenta en los accesos peatonales del TESCHA y se agudiza
en el turno vespertino, durante los 15 a 20 minutos previos al inicio de la
primera clase, cuando la mayor parte de los estudiantes llega en un intervalo
corto.

\subsection{La Importancia}
La importancia de este problema se aprecia en quiénes lo padecen:
\begin{enumerate}
	\item A los \textbf{estudiantes}, que pierden tiempo en las filas de
	      entrada, se exponen a llegar tarde a su primera clase y conviven en un
	      plantel donde el ingreso de personas ajenas no se controla con rigor.
	\item Al \textbf{personal de vigilancia}, que debe revisar cientos de
	      credenciales en pocos minutos sin ninguna herramienta que le ayude a
	      detectar una credencial ajena, vencida o apócrifa.
	\item A los \textbf{directivos y al Departamento de Control Escolar}, que no
	      disponen de datos de ingreso para detectar a estudiantes que dejan de acudir
	      al plantel ni para planear la apertura de accesos y el personal necesario en
	      cada horario.
\end{enumerate}

Y también en lo que provoca:
\begin{itemize}
	\item \textbf{Filas y tiempos muertos} al inicio de la jornada, que
	      provocan retardos en la primera clase. La teoría de colas, que ya se ha
	      propuesto para reducir los tiempos de espera de estudiantes en el área de
	      control escolar de otro plantel del Tecnológico Nacional de México
	      \parencite{talavera2025espera}, permite analizar este efecto a partir de la
	      tasa de llegadas y del tiempo de atención.
	\item \textbf{Riesgo de seguridad}, porque una credencial prestada, vencida
	      o ajena puede permitir el ingreso de personas que no pertenecen a la
	      comunidad del TESCHA.
	\item \textbf{Ausencia de datos}: sin un registro digital no es posible
	      saber cuántas personas ingresan ni en qué horarios, ni detectar a tiempo a
	      los estudiantes que dejan de acudir al plantel, una señal temprana de
	      posible deserción.
	\item \textbf{Carga operativa} para el personal de vigilancia, cuya
	      revisión depende únicamente de la atención de una persona en los momentos
	      de mayor afluencia.
\end{itemize}

\subsection{Lo que Desconocemos}
Para comprender mejor el problema se requiere información en cinco ejes.

\subsubsection{Situación Actual del Acceso en el TESCHA}
Es la información más importante, porque sin ella no es posible afirmar la
magnitud del problema. Se necesita conocer la matrícula por turno, el número
de accesos y de vigilantes, la tasa de llegada de estudiantes por minuto en el
horario pico, el tiempo que toma revisar a cada estudiante, la frecuencia con
que se olvidan o se prestan credenciales y la proporción de estudiantes que
cuenta con un celular propio. Estos datos se obtendrán con la encuesta del
Apéndice~A, la entrevista del Apéndice~B y la hoja de observación del
Apéndice~C.

\subsubsection{Arquitectura de Hardware y Tiempos de Respuesta}
\begin{itemize}
	\item \textbf{Latencia de lectura:} cuántos milisegundos toma al ESP32-CAM
	      capturar y decodificar el código, y cuántos más consume la consulta al
	      servidor por Wi-Fi. La suma debe ser menor que el tiempo de revisión manual
	      medido en el Apéndice~C para que el sistema tenga sentido. No se encontró
	      literatura arbitrada que reporte este tiempo en un ESP32-CAM, por lo que se
	      medirá de forma experimental con el prototipo; la referencia más cercana
	      disponible es un preprint que documenta el desempeño general de transmisión
	      de video del ESP32-CAM \parencite{nowroz2025esp32cam}, aunque no aborda la
	      decodificación de códigos QR.
	\item \textbf{Condiciones de lectura:} cómo afectan la iluminación, el
	      brillo de la pantalla del celular y la distancia al lector la tasa de
	      lecturas fallidas.
	\item \textbf{Durabilidad mecánica:} cuántos ciclos diarios deberá soportar
	      el mecanismo del torniquete según el flujo observado.
\end{itemize}

\subsubsection{Seguridad y Protección de Datos Personales}
\begin{itemize}
	\item \textbf{Validez del código:} qué duración del token equilibra la
	      seguridad frente a capturas de pantalla con la tolerancia a desfases de
	      reloj entre el celular y el servidor.
	\item \textbf{Comunicación segura:} uso de HTTPS entre el lector, el
	      servidor y el panel web para que los registros no puedan interceptarse ni
	      modificarse.
	\item \textbf{Marco legal:} como organismo público del Estado de México, el
	      TESCHA es sujeto obligado de la Ley de Protección de Datos Personales en
	      Posesión de Sujetos Obligados del Estado de México y Municipios
	      \parencite{edomex2017datos}, además de la Ley General de la materia
	      expedida en 2025 \parencite{lgpdppso2025}. Hay que identificar las
	      obligaciones que ambas imponen al sistema, como el aviso de privacidad y el
	      periodo de conservación de los registros de ingreso.
\end{itemize}

\subsubsection{Modelado de Datos e Indicadores}
\begin{itemize}
	\item \textbf{Estructura relacional:} diseño normalizado de las tablas de
	      estudiantes, accesos, dispositivos lectores y registros de ingreso en MySQL
	      o PostgreSQL.
	\item \textbf{Indicadores útiles:} qué reportes necesitan realmente
	      vigilancia y Control Escolar, por ejemplo, ingresos por franja horaria,
	      intentos rechazados y estudiantes sin ingresos durante varios días
	      consecutivos.
\end{itemize}

\subsubsection{Aceptación y Experiencia de Usuario}
\begin{itemize}
	\item \textbf{Disposición a usar el sistema:} si los estudiantes estarían
	      dispuestos a mostrar un QR en su celular y qué preocupaciones tienen sobre
	      el registro de su hora de entrada.
	\item \textbf{Ergonomía del lector:} altura y ángulo de la cámara para que
	      estudiantes de distintas estaturas lean su código sin detener el paso.
\end{itemize}

\subsection{La Viabilidad}
El uso de códigos QR para el control de acceso físico ya se ha propuesto en
la literatura \parencite{kao2011physical}, y su bajo costo frente a otras
tecnologías (Tabla~\ref{tab:comparativa}) hace viable construir y probar un
prototipo con los recursos del equipo. Frente a NFC, una comparación en
transacciones móviles encontró que el QR se procesa con mayor rapidez, aunque
no incorpora cifrado propio \parencite{harnaningrum2024comparison}; esa
carencia se compensa con la firma HMAC descrita en Delimitación. Además, el
estudiante usa su propio celular, por lo que no hace falta emitir tarjetas.
Frente al reconocimiento facial, el QR no requiere procesamiento pesado ni
trata datos biométricos, cuya aceptación entre estudiantes universitarios
depende en buena medida de la confianza y de la seguridad que perciben
\parencite{rukhiran2023user}, y cuyo tratamiento implicaría, además, las
obligaciones legales adicionales señaladas en ``Seguridad y Protección de
Datos Personales''.

\begin{table}[tbp]
	\caption{Comparación de Tecnologías de Validación de Acceso}
	\label{tab:comparativa}
	\begin{tabular}{lccc}
		\toprule
		Criterio                            & Código QR                & NFC     & Reconocimiento facial \\
		\midrule
		Costo de infraestructura            & Bajo                     & Medio   & Alto                  \\
		Medio que porta el estudiante       & Celular                  & Tarjeta & Ninguno               \\
		Complejidad de cómputo en el lector & Baja                     & Baja    & Alta                  \\
		Trata datos biométricos             & No                       & No      & Sí                    \\
		Riesgo de copia del identificador   & Medio\textsuperscript{a} & Bajo    & Bajo                  \\
		\bottomrule
	\end{tabular}
	\tablenote{Elaboración propia con base en
		\textcite{harnaningrum2024comparison} y \textcite{rukhiran2023user}. El costo
		y la complejidad son una valoración cualitativa del equipo.
		\textsuperscript{a}El riesgo corresponde a un QR estático; con un código
		dinámico firmado con HMAC se reduce considerablemente
		\parencite{subpratatsavee2015hmac}.}
\end{table}

Queda por conocer la viabilidad tecnológica desde el lado del estudiante, es
decir, si cuenta con un celular propio, con conexión a internet y con batería
suficiente al llegar al plantel; esto se indaga en la Sección~IV del
Apéndice~A.

La investigación es necesaria porque hoy no existen datos que midan el
problema en el TESCHA: se desconoce cuánto tiempo espera un estudiante para
entrar, con qué frecuencia se usan credenciales ajenas y si la comunidad
aceptaría un sistema basado en su celular. Levantar esa información y
contrastarla con el desempeño de un prototipo permitirá saber, con evidencia y
no con suposiciones, si esta solución de bajo costo ayuda a aliviar los
estragos del sistema actual y si justifica su adopción en el resto de los
accesos del plantel.


\section{Pregunta de Investigación}
A partir del problema descrito, la pregunta que guía esta investigación es la
siguiente: ¿en qué medida un sistema de acceso con código QR dinámico, leído
por un ESP32-CAM, reduce el tiempo promedio de ingreso por estudiante
respecto a la revisión visual actual, y en qué medida disminuye los ingresos
con credenciales ajenas? Con ello se busca saber si el sistema ayuda a
aliviar los estragos del control de acceso actual del TESCHA. Para
responderla se medirán tres indicadores, antes del piloto y durante él, con
la hoja del Apéndice~C: el tiempo promedio de validación por estudiante, la
longitud máxima de la fila en el horario pico y el número de incidencias de
identidad detectadas.


\section{¿Cómo Debemos Investigarlo?}
De acuerdo con los criterios revisados en clase, el tipo de investigación se
elige a partir del problema y de la pregunta planteados en las secciones
anteriores. Cada decisión se declara y se justifica enseguida.

\subsection{Finalidad}
La investigación es \textbf{aplicada}. Se entiende por aplicada la
investigación dirigida primordialmente a un fin práctico específico, en
contraste con la investigación pura o básica, cuyo propósito es adquirir
conocimiento sin un objetivo de aplicación predeterminado
\parencite{oecd2015frascati}; su valor está en producir evidencia para
resolver problemas concretos de la práctica \parencite{vargascordero2009aplicada}.
Este proyecto se ubica en el primer caso: no se pretende generar conocimiento
teórico nuevo sobre los códigos QR ni sobre la teoría de colas, temas que ya
cuentan con literatura consolidada, sino emplear ese conocimiento existente
para resolver un problema situado, el control de acceso de un plantel
determinado. El propósito del estudio es producir una decisión práctica —si
conviene o no adoptar el sistema en el resto de los accesos— y un producto
utilizable, el prototipo del Apéndice~D. Sus resultados valen para el caso
estudiado y no se presentan como una generalización.

\subsection{Enfoque}
El enfoque es \textbf{mixto, con predominancia cuantitativa}
\parencite{hernandezsampieri2014metodologia}, porque el problema exige dos
clases distintas de información. Los indicadores de tiempo
de validación, longitud de fila e incidencias de identidad (encuesta del
Apéndice~A y hoja de observación del Apéndice~C) son magnitudes medibles y se
analizan de forma cuantitativa, pues la pregunta de investigación pide saber
\emph{en qué medida} cambian. La entrevista del Apéndice~B aporta información
cualitativa que ninguna medición recoge: las razones, los criterios y las
dificultades que el personal de vigilancia y de Control Escolar expresa con
sus propias palabras.

\subsection{Alcance}
El alcance es \textbf{descriptivo y explicativo}. Es descriptivo en su primera
etapa, de carácter diagnóstico, porque el nivel de conocimiento disponible
sobre el acceso del TESCHA es nulo y hace falta documentar la situación actual
—cuánto se espera, cuántos entran, con qué frecuencia falla la
identificación— para establecer una línea base. Es explicativo en su segunda
etapa, porque una vez instalado el prototipo se busca establecer si el cambio
de mecanismo de validación explica la reducción observada en el tiempo de
ingreso y en las incidencias de identidad, y no solo constatar que ocurrió.

\subsection{Obtención de Información}
La obtención de información es \textbf{mixta}, documental y de campo. Es
documental en la fundamentación: la elección del código QR frente a NFC y al
reconocimiento facial, el esquema de firma HMAC y el diseño cuasiexperimental
se sustentan en fuentes ya publicadas, revisadas en las secciones anteriores.
Es de campo en la parte central del estudio, porque los datos que dan
respuesta a la pregunta de investigación no existen en ninguna fuente
documental: no hay registro alguno del flujo de acceso del TESCHA, de modo
que deben levantarse en el sitio, en el acceso piloto y en el horario en que
el problema se presenta.

\subsection{Técnica e Instrumento de Recolección de Datos}
Se emplean tres pares de técnica e instrumento, cada uno elegido por el tipo
de información que se necesita obtener:
\begin{itemize}
	\item \textbf{Encuesta}, con un \textbf{cuestionario} de preguntas cerradas
	      y escala de Likert (Apéndice~A). Se elige porque la percepción de espera, la
	      frecuencia de olvido de credencial, la disposición a usar el sistema y la
	      viabilidad tecnológica del estudiante son datos que solo el propio
	      estudiante puede reportar, y se requieren de una población numerosa en poco
	      tiempo, con respuestas comparables entre sí y de forma anónima.
	\item \textbf{Entrevista semiestructurada}, con una \textbf{guía de
		      entrevista} (Apéndice~B). Se elige porque el personal de vigilancia y de
	      Control Escolar son informantes clave, pocos en número, que poseen
	      información que ningún cuestionario anticipa: el procedimiento real de
	      revisión, los casos excepcionales y las condiciones institucionales para
	      tratar datos de los estudiantes. El formato semiestructurado permite
	      repreguntar sobre lo que aparezca.
	\item \textbf{Observación estructurada no participante}, con una \textbf{hoja
		      de registro} (Apéndice~C). Se elige porque el tiempo de revisión y la
	      longitud de la fila son hechos observables que deben medirse
	      directamente con cronómetro y conteo; preguntarlos produciría estimaciones
	      subjetivas, no mediciones. Es además el único instrumento que se aplica dos
	      veces con el mismo formato, lo que hace comparables la línea base y el
	      piloto.
\end{itemize}

\subsection{Control de Variables}
En cuanto al control de variables, la investigación es \textbf{experimental},
aunque no en sentido estricto, sino \textbf{cuasiexperimental}. Sí hay
manipulación deliberada de la variable independiente —el mecanismo de
validación de identidad, que pasa de la revisión visual de credenciales al
escaneo del QR dinámico— y se observa su efecto sobre las variables
dependientes: el tiempo promedio de validación por estudiante, la longitud
máxima de la fila y el número de incidencias de identidad. Lo que no es
posible es el control propio de un experimento puro: no hay asignación
aleatoria de los estudiantes ni grupo de control, porque el acceso piloto
atiende a quien llega. Se mantienen constantes las condiciones que sí están al
alcance del equipo —el mismo acceso, el mismo turno, la misma franja horaria,
los mismos observadores y el mismo instrumento en ambas mediciones— y se
registra el número de puntos de revisión activos ($c$) para no atribuir al
prototipo una mejora que provenga de haber abierto más filas.

\subsection{Temporalidad}
Según el criterio temporal, que distingue entre recolectar datos en un solo
momento y recolectarlos en varios puntos en el tiempo para hacer inferencias
sobre el cambio \parencite{hernandezsampieri2014metodologia}, esta
investigación es \textbf{longitudinal}, con mediciones repetidas. El fenómeno
no se mide una sola vez, sino en dos series separadas por la intervención:
cinco días hábiles consecutivos de observación antes del piloto y una segunda
serie equivalente durante él, sobre el mismo acceso y con la misma hoja de
registro. Esta temporalidad la exige la pregunta de investigación, que pide
comparar un antes y un después; un estudio transversal, de una sola medición,
permitiría describir la situación actual pero no atribuir ningún cambio al
prototipo.

Conviene precisar el uso del término, porque el eje
transversal/longitudinal suele desarrollarse para los diseños no
experimentales, en los que el investigador observa sin administrar ningún
estímulo, y sus subtipos —como el diseño de panel— se definen en ese marco
\parencite{hernandezsampieri2014metodologia}. Aquí sí hay un estímulo
administrado por el equipo, el prototipo, de modo que el diseño pertenece a la
familia cuasiexperimental descrita más adelante. Por eso «longitudinal» se
emplea en este trabajo en su sentido temporal —varias mediciones a lo largo
del tiempo sobre las mismas unidades de observación— y no como la etiqueta de
un diseño no experimental de panel.

\subsection{Diseño}
En consecuencia, el diseño es de campo y cuasiexperimental, del tipo
preprueba/posprueba con un solo grupo —el acceso piloto— y sin grupo de
control \parencite{campbellstanley1963, hernandezsampieri2014metodologia},
dado que no es posible operar dos accesos idénticos en paralelo bajo las
mismas condiciones. Un diseño equivalente se ha empleado recientemente para
evaluar, también sin grupo de control, un piloto de intervención basado en
código QR dentro de una institución \parencite{ortiz2026qrpilot}. Entre las
amenazas a la validez interna de este diseño destacan la historia y la
maduración \parencite{cranmer2017pretestposttest} —por ejemplo, un cambio de
horario escolar o la simple costumbre de los estudiantes podrían alterar los
tiempos sin que el prototipo tenga parte en ello—, por lo que los resultados
se interpretarán como evidencia orientativa y no concluyente. Los tres
instrumentos se aplican antes del piloto para obtener la línea base y de
nuevo durante él para comparar los resultados, como se detalla en el Plan de
Acción Preliminar (Apéndice~D).


\printbibliography

% ----------------------------- Apéndices -------------------------------
\appendix

\clearpage
\section{Encuesta a Estudiantes sobre el Acceso al Plantel}
\apendiceindice{A}{Encuesta a Estudiantes sobre el Acceso al Plantel}

\subsection{Aplicación}
La encuesta se dirige a estudiantes de todas las carreras y turnos del TESCHA,
es anónima y toma alrededor de cinco minutos. Puede aplicarse impresa o
trasladarse a un formulario en línea. El tamaño de muestra se calcula con la
fórmula para proporciones en una población finita:
\[
	n = \frac{N z^{2} p q}{e^{2}(N-1) + z^{2} p q}
\]
donde $N$ es la matrícula total del plantel (dato que proporcionará Control
Escolar), $z = 1.96$ para un nivel de confianza del 95\,\%, $p = q = 0.5$ y
$e$ es el margen de error. Para una matrícula grande, $n$ tiende a 384
estudiantes con $e = 0.05$ y a 196 con $e = 0.07$.

\subsection{Instrumento}
\noindent\textit{Aviso de privacidad simplificado.} Esta encuesta forma parte de
un proyecto académico de Ingeniería en Sistemas Computacionales del TESCHA. No
se solicita tu nombre ni tu matrícula, y las respuestas se usarán únicamente de
forma agregada con fines académicos. Responder es voluntario.

\medskip
\noindent\textbf{Sección I. Datos generales}
\begin{enumerate}
	\item Carrera: \linea[6cm] \qquad Semestre: \linea[1.5cm]
	\item Turno: \op Matutino \quad \op Vespertino
	\item ¿Por cuál acceso entras normalmente? \linea[6cm]
	\item ¿Con cuánta anticipación sueles llegar a tu primera clase?\\
	      \op Más de 20 min \quad \op De 10 a 20 min \quad \op Menos de 10 min
	      \quad \op Llego después de que inicia
\end{enumerate}

\noindent\textbf{Sección II. Experiencia con el acceso actual}
\begin{enumerate}\setcounter{enumi}{4}
	\item En un día normal, ¿cuánto tiempo esperas en la fila para entrar?\\
	      \op No hay fila \quad \op Menos de 1 min \quad \op De 1 a 3 min
	      \quad \op De 3 a 5 min \quad \op Más de 5 min
	\item En este semestre, ¿cuántas veces has llegado tarde a clase por la fila
	      de entrada?\\
	      \op Ninguna \quad \op 1 o 2 \quad \op De 3 a 5 \quad \op Más de 5
	\item En este semestre, ¿cuántas veces has olvidado o perdido tu credencial?\\
	      \op Ninguna \quad \op 1 o 2 \quad \op 3 o más
	\item Cuando alguien llega sin credencial, ¿qué ocurre normalmente?\\
	      \op Pasa sin revisión \quad \op Se registra de forma manual
	      \quad \op No lo dejan pasar \quad \op No lo sé
	\item ¿Has visto entrar a alguien con la credencial de otra persona o a
	      alguien que no es estudiante?\\
	      \op Nunca \quad \op Una vez \quad \op Varias veces
	      \quad \op Prefiero no responder
\end{enumerate}

\noindent\textbf{Sección III. Opinión.} Marca qué tan de acuerdo estás con cada
afirmación (1 = totalmente en desacuerdo, 5 = totalmente de acuerdo).

\medskip
\noindent
\begin{tabular}{p{0.62\textwidth}ccccc}
	\toprule
	Afirmación                                                               & 1   & 2   & 3   & 4   & 5   \\
	\midrule
	10. La revisión actual en la entrada es rápida.                          & \op & \op & \op & \op & \op \\
	11. La revisión actual impide que entren personas ajenas al TESCHA.      & \op & \op & \op & \op & \op \\
	12. Me siento con seguridad dentro del plantel.                          & \op & \op & \op & \op & \op \\
	13. Usaría un código QR en mi celular para entrar al plantel.            & \op & \op & \op & \op & \op \\
	14. Me incomoda que el sistema registre la hora a la que entro.          & \op & \op & \op & \op & \op \\
	15. Prefiero un código QR antes que un sistema de reconocimiento facial. & \op & \op & \op & \op & \op \\
	\bottomrule
\end{tabular}

\medskip
\noindent\textbf{Sección IV. Viabilidad tecnológica}
\begin{enumerate}\setcounter{enumi}{15}
	\item ¿Llevas contigo todos los días un teléfono inteligente propio?\\
	      \op Sí \quad \op Casi siempre \quad \op No
	\item Al llegar al plantel, ¿tu celular tiene conexión a internet (datos o
	      Wi-Fi)?\\
	      \op Siempre \quad \op A veces \quad \op Nunca
	\item Al llegar al plantel, ¿tu celular suele tener batería suficiente?\\
	      \op Siempre \quad \op Casi siempre \quad \op A veces \quad \op Casi nunca
	\item Si un día no pudieras usar tu celular, ¿qué alternativa preferirías?\\
	      \op Registro manual con vigilancia \quad \op Credencial con QR impreso
	      \quad \op Otra: \linea[3cm]
	\item Comentarios o sugerencias sobre el acceso al plantel:\\[0.4em]
	      \linea[\linewidth]\\[0.8em]\linea[\linewidth]
\end{enumerate}


\clearpage
\section{Guía de Entrevista al Personal de Vigilancia y de Control Escolar}
\apendiceindice{B}{Guía de Entrevista al Personal de Vigilancia y de Control Escolar}
Entrevista semiestructurada de 10 a 15 minutos, con al menos un integrante del
personal de vigilancia del acceso piloto y uno del Departamento de Control
Escolar. Se solicitará consentimiento verbal y no se registrarán nombres.

\subsection{Personal de Vigilancia}
\begin{enumerate}
	\item ¿Cómo se revisa hoy la identidad de quien entra? ¿Qué se verifica de
	      la credencial?
	\item ¿Cuántas personas atienden el acceso en el horario de mayor afluencia y
	      a qué hora ocurre?
	\item ¿Qué se hace cuando un estudiante llega sin credencial?
	\item ¿Con qué frecuencia detectan credenciales ajenas, vencidas o personas
	      que no son estudiantes?
	\item ¿Qué les facilitaría el trabajo en la entrada?
\end{enumerate}

\subsection{Departamento de Control Escolar}
\begin{enumerate}
	\item ¿Cuál es la matrícula actual por turno?
	\item ¿Existe hoy algún registro de ingresos al plantel? ¿Quién lo consulta?
	\item ¿Qué información sobre el ingreso de los estudiantes les sería útil
	      para detectar ausentismo o riesgo de deserción?
	\item ¿Qué condiciones institucionales habría que cumplir para usar datos de
	      los estudiantes en un sistema de acceso (aviso de privacidad, autorización,
	      resguardo)?
\end{enumerate}


\clearpage
\section{Hoja de Observación del Flujo en el Acceso}
\apendiceindice{C}{Hoja de Observación del Flujo en el Acceso}

\subsection{Procedimiento}
La observación se realiza en el acceso piloto durante cinco días hábiles
consecutivos, desde 30 minutos antes hasta 10 minutos después del inicio de la
primera clase del turno vespertino, en intervalos de cinco minutos. Participan
dos observadores: el primero cuenta los ingresos y anota la longitud máxima de
la fila en cada intervalo; el segundo cronometra, para cinco estudiantes
elegidos al azar por intervalo, el tiempo desde que llegan al punto de revisión
hasta que lo cruzan. La misma hoja se utilizará durante el piloto para comparar
los resultados.

\subsection{Formato}
\noindent Fecha: \linea[3cm] \quad Día: \linea[2.5cm] \quad Acceso:
\linea[3cm]\\[0.5em]
Hora de inicio de clases: \linea[2cm] \quad Puntos de revisión activos ($c$):
\linea[1.5cm] \quad Observadores: \linea[3cm]

\medskip
\noindent
{\small
	\begin{tabular}{|p{2.2cm}|p{1.3cm}|p{1.3cm}|p{4.2cm}|p{1.2cm}|p{1.2cm}|p{1.2cm}|}
		\hline
		\textbf{Intervalo (min)}         & \textbf{Ingresos}  & \textbf{Fila máx.} &
		\textbf{Tiempos de revisión (s)} & \textbf{Sin cred.} & \textbf{Sin
		revisión}                        & \textbf{Rechaz.}                                    \\
		\hline
		$-30$ a $-25$                    &                    &                    &   &  &  & \\[0.8em] \hline
		$-25$ a $-20$                    &                    &                    &   &  &  & \\[0.8em] \hline
		$-20$ a $-15$                    &                    &                    &   &  &  & \\[0.8em] \hline
		$-15$ a $-10$                    &                    &                    &   &  &  & \\[0.8em] \hline
		$-10$ a $-5$                     &                    &                    &   &  &  & \\[0.8em] \hline
		$-5$ a $0$                       &                    &                    &   &  &  & \\[0.8em] \hline
		$0$ a $+5$                       &                    &                    &   &  &  & \\[0.8em] \hline
		$+5$ a $+10$                     &                    &                    &   &  &  & \\[0.8em] \hline
	\end{tabular}}

\subsection{Análisis}
Con los datos de cada intervalo se calculan la tasa de llegadas
$\lambda = \mathrm{ingresos}/5$ (estudiantes por minuto), el tiempo medio de
revisión $\bar{t}$ en segundos, la tasa de servicio de un punto de revisión
$\mu = 60/\bar{t}$ y la utilización del acceso:
\[
	\rho = \frac{\lambda}{c\,\mu}
\]
Cuando $\rho \geq 1$, llegan más estudiantes de los que el acceso puede atender
y la fila crece durante ese intervalo. Comparar $\bar{t}$, $\rho$ y la fila
máxima antes del piloto y durante él permite comprobar si el prototipo reduce
el tiempo de ingreso.


\clearpage
\section{Plan de Acción Preliminar}
\apendiceindice{D}{Plan de Acción Preliminar}
\begin{enumerate}
	\item \textbf{Diagnóstico:} aplicar la encuesta, las entrevistas y la
	      observación de los Apéndices A a C para obtener la línea base del acceso
	      piloto.
	\item \textbf{Diseño de hardware y firmware:} integrar el ESP32-CAM con el
	      mecanismo del torniquete y programar la lectura del código y la consulta al
	      servidor.
	\item \textbf{Desarrollo del servidor:} implementar la generación y
	      validación de QR dinámicos firmados con HMAC, la base de datos de ingresos y
	      el panel web con los indicadores definidos.
	\item \textbf{Piloto controlado:} instalar el prototipo en el acceso piloto
	      durante un periodo acotado, sin retirar la revisión manual.
	\item \textbf{Evaluación:} repetir la observación del Apéndice~C, comparar
	      los resultados con la línea base y documentar los ajustes necesarios.
\end{enumerate}


\end{document}
